\section{Matrix product approximation of a frustration-free ground vector}

The frustration-free argument of \cite[Section~5]{Arad2013AreaLaw} filters a
state of constant overlap once more, with a projector of degree $s^2$ on a
window whose width grows like the cube root of the squared logarithm of the
accuracy. The exponent of the gap in the shrinking factor is $+1/2$; the
printed bond dimension omits the factor $\log\ell$ of the polynomial rank
estimate, which is retained here \cite{gap:aklv13_frustration_free_mps_rate}.

\begin{lemma}[A filtered product vector]
    \label{lem:aklv_ff_product_filtered_output}
    \lean{TNLean.MPS.AreaLaw.exists_cut_product_filtered_output}
    \leanok
    Let $K$ fix a unit vector $\phi$ of an $L$-site chain, preserve
    $\phi^\perp$, and satisfy $\|Kx\|^2\le\delta$ for every unit
    $x\perp\phi$, with $\delta\ge0$. Suppose
    $\operatorname{OSR}_k(K)\le D$ with $D\ge1$ and $D\delta\le1/2$.
    Then there are a unit vector $\chi$, which is a product across the cut
    $k$, with $|\langle\phi,\chi\rangle|\ge1/\sqrt{2D}$, and a phase $c$,
    such that $K\chi\ne0$ and $\psi=cK\chi/\|K\chi\|$ satisfies
    \begin{align*}
        \operatorname{SR}_k(\psi) & \le D,       &
        \|\psi-\phi\|^2           & \le4D\delta.
    \end{align*}
    This is \cite[Lemma~2.2]{Arad2013AreaLaw} followed by one application
    of the filter.
\end{lemma}
\begin{proof}
    \leanok
    Transport $K$ and $\phi$ along the identification of the chain with
    the two sides of the cut. The maximal product overlap is at least
    $1/\sqrt{2D}$. The phase-aligned normalized image has error at most
    $2\delta\|\chi-\langle\phi,\chi\rangle\phi\|^2/|\langle\phi,\chi\rangle|^2
        \le4D\delta$, and its Schmidt rank is at most $D$ because $\chi$ is
    a product.
\end{proof}

\begin{lemma}[Schmidt rank of a filtered vector]
    \label{lem:aklv_normalized_cut_schmidt_rank}
    \lean{TNLean.MPS.AreaLaw.normalized_cut_schmidtRank_le}
    \leanok
    If $c\ne0$ and $Kx\ne0$, then
    $\operatorname{SR}_k(cKx/\|Kx\|)\le\operatorname{OSR}_k(K)\operatorname{SR}_k(x)$.
\end{lemma}
\begin{proof}
    \leanok
    Nonzero scalars do not change the Schmidt rank, and an operator of
    operator Schmidt rank $D$ multiplies the Schmidt rank by at most $D$.
\end{proof}

\begin{definition}[Cube-root accuracy parameters]
    \label{def:aklv_ff_accuracy_parameters}
    \lean{TNLean.MPS.AreaLaw.aklvFrustrationFreeIncrement,
        TNLean.MPS.AreaLaw.aklvFrustrationFreeRankBound}
    \leanok
    \uses{def:aklv_accuracy_parameters}
    For $0<g\le1$, squared accuracy $0<\rho\le1$, local dimension $q$ and
    initial width $s_0$, put $z=\log(128/\rho)$ and
    \begin{align*}
        r & =\left(\frac{z^2}{g}\right)^{1/3}, &
        m & =\lceil6r\rceil,                   &
        s & =s_0+2m,
    \end{align*}
    and define the rank budget
    \begin{align*}
        B=\max\Bigl(\bigl(q(s_0^2+1)\bigr)^{32(2s_0+1)}
        \bigl(q(s^2+1)\bigr)^{32(2s+1)},\ q^{s-\lfloor s/2\rfloor}\Bigr).
    \end{align*}
\end{definition}

\begin{lemma}[The cube-root scale]
    \label{lem:aklv_ff_increment_bounds}
    \lean{TNLean.MPS.AreaLaw.aklvFrustrationFreeIncrement_bounds}
    \leanok
    \uses{def:aklv_ff_accuracy_parameters}
    For $0<g\le1$ and $0<\rho\le1$ one has
    $r\ge1$, $r^3=z^2/g$, $6r\le m\le7r$ and $r\le z/g$.
\end{lemma}
\begin{proof}
    \leanok
    \uses{thm:aklv_accuracy_parameter_bounds}
    Since $z\ge1$ and $g\le1$, $z^2/g\ge1$, so its cube root is at least one.
    The ceiling gives $6r\le m<6r+1\le7r$. Finally
    $r^3=z^2/g\le z^3/g^3$.
\end{proof}

\begin{lemma}[Shrinking of the projector of degree $s^2$]
    \label{lem:aklv_ff_accuracy_shrink}
    \lean{TNLean.MPS.AreaLaw.aklvFrustrationFree_shrink_le}
    \leanok
    Let $0<g$, $0<\rho\le1$, $t=1024/g$ and $z=\log(128/\rho)$. If an
    integer $s$ satisfies $s\ge1/g$ and $gs^3\ge1728z^2$, then
    \begin{align*}
        2\exp\left(-2s^2\sqrt{\frac{g/2}{2t+s}}\right)\le\frac{\rho}{64}.
    \end{align*}
    The exponent of the gap is the one of the proof of
    \cite[Theorem~4.3]{Arad2013AreaLaw}.
\end{lemma}
\begin{proof}
    \leanok
    \uses{thm:aklv_accuracy_parameter_bounds}
    Since $s\ge1/g$, $2t+s\le2049s$. Hence the square of the exponent is at
    least $4s^3g/4098\ge z^2$. The exponential is therefore at most
    $e^{-z}=\rho/128$.
\end{proof}

\begin{lemma}[The cube-root logarithmic rank bound]
    \label{lem:aklv_ff_rank_log}
    \lean{TNLean.MPS.AreaLaw.aklvFrustrationFreeRankBound_log_le}
    \leanok
    \uses{def:aklv_ff_accuracy_parameters}
    Let $q\ge2$, $0<g\le1$, $0<\rho\le1$, $s_0\ge1$, $R,W\ge0$, and put
    $u=1+\log q+\log(1/g)$ and $D_0=(q(s_0^2+1))^{32(2s_0+1)}$. If
    $\log D_0\le Ru^3/g$ and $s_0\le Wu^2/g$, then
    \begin{align*}
        \log B\le(10001R+W)\frac{u^3}{g}
        +8400000\,r\bigl(u+\log(z+1)\bigr).
    \end{align*}
\end{lemma}
\begin{proof}
    \leanok
    \uses{lem:aklv_ff_increment_bounds,thm:aklv_quadratic_degree_log_rank}
    Put $D_1=(q(s^2+1))^{32(2s+1)}$, the rank formula at degree $s^2$.
    If $m\le s_0$, then $s\le3s_0$, so $s^2\le66s_0^2$ and
    $\log D_1\le10000\log D_0$.
    If $m>s_0$, then $s\le3m$, so $s^2\le66m^2$, while $m\le7r\le7z/g$.
    The second quadratic-degree estimate at the accuracy $7z/g$ gives
    $\log D_1\le400000\,m(1+\log q+\log(7z/g+1))$, and the last factor is at
    most $3(u+\log(z+1))$. The boundary rank costs at most
    $(s_0+2m)\log q\le Wu^3/g+14r u$.
\end{proof}

\begin{lemma}[Normalization of a frustration-free gap]
    \label{lem:aklv_ff_dimension_gap}
    \lean{TNLean.MPS.AreaLaw.aklv_frustrationFree_dimension_and_gap_bounds}
    \leanok
    Let $0\le h_j\le I$ be nearest-neighbour terms of an $L$-site chain, and
    $\Omega$ a unit vector with $H\Omega=0$ and
    $H\ge g(I-|\Omega\rangle\langle\Omega|)$, where $H=\sum_jh_j$ and $g>0$.
    If $g$ is an eigenvalue of $H$, then $q\ge2$ and $g\le1$.
\end{lemma}
\begin{proof}
    \leanok
    \uses{thm:chain_positive_nearest_neighbor_gap_upper}
    The two distinct eigenvalues $0$ and $g$ force $q\ge2$ and $L\ge1$.
    The endpoint trial vector then gives $g\le1$.
\end{proof}

\begin{lemma}[Two filters at one cut]
    \label{lem:aklv_ff_two_filter_approximant}
    \lean{TNLean.MPS.AreaLaw.exists_aklv_frustrationFree_two_filter_approximant}
    \leanok
    Let $q\ge1$, $0<g\le1$, $t=1024/g$, and let $\Omega$ be a unit vector with
    $H\Omega=0$ and $H\ge g(I-|\Omega\rangle\langle\Omega|)$ for a sum $H$ of
    nearest-neighbour terms $0\le h_j\le I$ on an $L$-site chain. Let
    $s_0\ge6$, $s=s_0+2m$ and $a+s+1\le L$, and suppose
    $D_0\delta_0\le10^{-5}$, where $D_0=(q(s_0^2+1))^{32(2s_0+1)}$ and
    $\delta_0=4\exp(-4s_0^2\sqrt{(g/2)/(2t+s_0)})$. Put
    $b=2\exp(-2s^2\sqrt{(g/2)/(2t+s)})$ and
    $D_1=(q(s^2+1))^{32(2s+1)}$. Then at the cut $k=a+\lfloor s/2\rfloor$
    there is a unit vector $y$ with
    \begin{align*}
        \|y-\Omega\|^2         & \le8b^2,    &
        \operatorname{SR}_k(y) & \le D_0D_1.
    \end{align*}
\end{lemma}
\begin{proof}
    \leanok
    \uses{lem:aklv_frustration_free_filter,
        lem:aklv_ff_product_filtered_output,
        lem:aklv_normalized_cut_schmidt_rank}
    The window of width $s_0$ with left end $a+m$ has the same centre $k$.
    Its filter of degree $s_0^2$ fixes $\Omega$, and
    Lemma~\ref{lem:aklv_ff_product_filtered_output} gives a vector $x$ of rank
    at most $D_0$ with $\|x-\Omega\|^2\le4D_0\delta_0\le1/16$, hence
    $|\langle\Omega,x\rangle|\ge1/2$. The filter of degree $s^2$ on the
    window of width $s$ has rank at most $D_1$ and shrinking amplitude $b$.
    Its phase-aligned normalized image $y$ of $x$ satisfies
    $\|y-\Omega\|^2\le2b^2\|x\|^2/|\langle\Omega,x\rangle|^2\le8b^2$, and its
    rank is at most $D_0D_1$ by Lemma~\ref{lem:aklv_normalized_cut_schmidt_rank}.
\end{proof}

\begin{lemma}[Frustration-free approximants at every cut]
    \label{lem:aklv_ff_per_cut}
    \lean{TNLean.MPS.AreaLaw.exists_aklv_frustrationFree_per_cut_approximation}
    \leanok
    There is a universal $C>0$ with the following property. Let $q\ge2$, let
    $0\le h_j\le I$ be nearest-neighbour terms on an $L$-site chain, and let
    $\Omega$ be a unit vector with $H\Omega=0$ and
    $H\ge g(I-|\Omega\rangle\langle\Omega|)$ for $0<g\le1$. For every
    $0<\rho\le1$ there is $D$ with
    \begin{align*}
        \log D\le C\left[\frac{u^3}{g}+\left(\frac{z^2}{g}\right)^{1/3}
            \bigl(u+\log(z+1)\bigr)\right],
    \end{align*}
    where $u=1+\log q+\log(1/g)$ and $z=\log(128/\rho)$, such that every cut
    $0\le k\le L$ has a unit vector $y$ with $\|y-\Omega\|^2\le\rho$ and
    $\operatorname{SR}_k(y)\le D$.
\end{lemma}
\begin{proof}
    \leanok
    \uses{lem:aklv_frustration_free_scalar_parameters,
        lem:aklv_ff_two_filter_approximant,
        lem:aklv_ff_accuracy_shrink,lem:aklv_ff_rank_log,
        thm:aklv_exact_boundary_rank}
    Take $s_0$, $t=1024/g$ and $D_0$ from the frustration-free scalar
    parameters, so $s_0\ge21$, $s_0\ge1/g$, $s_0\le Wu^2/g$,
    $\log D_0\le Ru^3/g$ and $D_0\delta_0\le10^{-5}$. Let $D=B$.
    If the window of width $s$ centred at $k$ fits,
    Lemma~\ref{lem:aklv_ff_two_filter_approximant} gives $y$ of rank at most
    $D_0D_1$ with $\|y-\Omega\|^2\le8b^2$. Since $s\ge12r$ and $s\ge s_0$,
    Lemma~\ref{lem:aklv_ff_accuracy_shrink} gives $b\le\rho/64$, so
    $8b^2\le\rho$. If the window does not fit, take $y=\Omega$, of rank at
    most $q^{s-\lfloor s/2\rfloor}$.
    Lemma~\ref{lem:aklv_ff_rank_log} gives the logarithmic bound with
    $C=10001R+W+8400000$.
\end{proof}

\begin{theorem}[Frustration-free matrix product approximation]
    \label{thm:aklv_frustration_free_mps}
    \lean{TNLean.MPS.AreaLaw.exists_aklv_frustrationFree_sequential_mps_approximation}
    \leanok
    \uses{def:aklv_sequential_spectral_approximant}
    There is a universal $C>0$ with the following property. Let $0\le h_j\le I$
    be nearest-neighbour terms on an $L$-site chain of local dimension $q$,
    and let $\Omega$ be a unit vector with $H\Omega=0$ and
    $H\ge g(I-|\Omega\rangle\langle\Omega|)$ for $H=\sum_jh_j$ and $g>0$, where
    $g$ is an eigenvalue of $H$. For $0<\epsilon\le1$ put
    \begin{align*}
        u & =1+\log q+\log\frac1g,             &
        z & =\log\frac{2048(L+1)}{\epsilon^2}.
    \end{align*}
    Then there is $D$ with
    \begin{align*}
        \log D\le C\left[\frac{u^3}{g}+\left(\frac{z^2}{g}\right)^{1/3}
            \bigl(u+\log(z+1)\bigr)\right]
    \end{align*}
    such that the normalized sequential truncation $\Phi$ of $\Omega$ at bond
    dimension $D$ is an open-boundary matrix product state with
    $\|\Phi-\Omega\|\le\epsilon$. For $\epsilon=L^{-c}$ and fixed $q$ and $g$
    this gives $\log D=O(\log^{2/3}L\,\log\log L)$. The printed bound
    $\exp(O(g^{-1/3}\log^{2/3}L))$ of \cite[Section~5]{Arad2013AreaLaw} is
    not asserted: it omits the factor $\log\ell$ of the polynomial rank
    estimate, as recorded in \cite{gap:aklv13_frustration_free_mps_rate}.
\end{theorem}
\begin{proof}
    \leanok
    \uses{lem:aklv_ff_per_cut,lem:aklv_ff_dimension_gap,
        thm:aklv_cut_approximation_tail,thm:aklv_sharp_sequential_mps}
    The actual gap gives $q\ge2$ and $g\le1$ by
    Lemma~\ref{lem:aklv_ff_dimension_gap}. Apply
    Lemma~\ref{lem:aklv_ff_per_cut} with $\rho=\epsilon^2/(16(L+1))$, so that
    $z=\log(128/\rho)$. At every cut the discarded mass of $\Omega$ beyond
    $D$ is at most $\rho$, by Lemma~\ref{thm:aklv_cut_approximation_tail}.
    The sequential truncation estimate then gives $\Phi$ with
    $\|\Phi-\Omega\|\le\epsilon$ and every internal cut rank at most $D$.
\end{proof}
